\scanpage{216}
\mspage{165}

Notons aussi qu'étant exact à gauche, le foncteur $\varphi_X$ de $X^\wedge$ dans $\mathrm{Cat}_{/\unsure{X|x}}$ transforme mono en mono. D'autre part :

\uline{Lemme 2} Les monomorphismes qui sont des isom. \deleted{\unsure{ouverts|ouvertes}} locaux sont exactement les imm. ouvertes. Donc \deleted{les} dans $\underline{\mathrm{Fib\,disc}}\,X$, les monomorphismes sont les immersions ouvertes, et le foncteur $\varphi_X\colon F\longmapsto X/F$ transforme un monomorphisme en \unsure{immersion ouverte|immersions ouvertes}.

Revenons alors au diagramme cocartésien $(\unsure{*|\alpha})$ (p. 161), et au diagramme correspondant dans Cat (et $=$ dans $\underline{\mathrm{Fib\,disc}}\,X$)

\diagram{\begin{tikzcd}[ampersand replacement=\&]
X/G\arrow[r,hook,"X/i"]\arrow[d,"X/p"']\& X/F\arrow[d,"X/q"] \& T\unsure{{:}|\prime|{}}\arrow[l,dotted,hook]\arrow[d,dotted]\\
X/\unsure{G'|G''}\arrow[r,hook,"X/i'"]\& X/F' \& T'\arrow[l,dotted,hook] .
\end{tikzcd}}

Comme $(*)$ est cartésien, le carré précédent est cartésien, par ce qui précède. Je veux prouver qu'il est également cocartésien\unsure{,|.}
% archive sheet number: 216
